\thispagestyle{empty}
{\sffamily\fontsize{9}{11}\selectfont\color{NEBGold}NEBIOT / PROJECT DOCUMENTATION\par}
\vspace{4mm}
\begin{center}\includegraphics[width=.94\linewidth]{assets/NEBIOT-English-6000px.png}\end{center}
{\sffamily\fontsize{8}{10}\selectfont\color{NEBGray}Program authority: Nuevo Estándar Biotropical — NEBIOT S.A.S.\par}
\vspace{5mm}
{\fontsize{32}{36}\selectfont\color{NEBGreen}\DocTitle\par}
\vspace{2mm}\textcolor{NEBGold}{\rule{\linewidth}{.6pt}}
\vspace{2mm}
{\sffamily\itshape\fontsize{11}{14}\selectfont\color{NEBGreen}Implementation, evidence and period-specific environmental accounts\par}
\vspace{2mm}
{\sffamily\fontsize{8.5}{11}\selectfont\color{NEBGray}Governing standard: NEBIOT 0.8.2 | Template edition 1.3\par}
\vspace{5mm}
{\fontsize{17}{21}\selectfont\itshape\color{NEBGreen}\ProjectName\par}
\vspace{3mm}
{\sffamily\fontsize{9}{12}\selectfont
\textbf{Project ID:} \ProjectID\par
\textbf{Country / jurisdiction:} \ProjectJurisdiction\par
\textbf{Project document version:} \ProjectVersion\par
\textbf{Issue date:} \ProjectDate\par\textbf{Monitoring interval: } \MonitoringStart{} -- \MonitoringEnd\par\textbf{Report number: } \MonitoringReportID\par}
\vspace{4mm}
\BlockTitle{Implementing organizations and stakeholders}
\expandafter\LogoSlot\expandafter{IMPLEMENTING ENTITY}{\ImplementerLogo}{\ImplementerName}\hfill
\LogoSlot{PROJECT PROPONENT}{\ProponentLogo}{\ProponentName}\hfill
\LogoSlot{RIGHTS HOLDERS / STAKEHOLDERS}{\StakeholderLogo}{\StakeholderName}\par
\vspace{2mm}
{\sffamily\fontsize{8}{10}\selectfont\color{NEBGray}Use authorized logos and identify each organization’s actual role. Logo display does not establish consent, project ownership, endorsement, validation or admission. Remove unused partner spaces from the completed submission.\par}
\vspace{2mm}
{\sffamily\fontsize{7.5}{9.5}\selectfont\color{NEBGray}The English logo wording is a translation, not a separate legal entity.\par}

\clearpage
{\fontsize{23}{27}\selectfont\color{NEBGreen}License and use of this template\par}
\vspace{4mm}
This document is issued by Nuevo Estándar Biotropical—NEBIOT S.A.S. Henry Arellano-Peña is the methodological author of NEB001 and NEB002. Rights in NEBIOT's protected materials remain with their respective rights holders. Commercial use, reproduction, adaptation, redistribution, and sublicensing of protected material are subject to the applicable license or specific written permission from the authorized rights holder; permissions for use within the NEBIOT program shall be recorded in FRM-07 under S-21. Access to this document does not itself grant a commercial license, transfer ownership, authorize sublicensing, confer program admission, or authorize certification or issuance of environmental units. Supporting publications and third-party data, software, and images retain their respective licenses and attribution requirements; existing lawful permissions and uses permitted by applicable law remain unaffected. Authorized independent reviewers shall receive the access necessary for inspection, recalculation, and reporting under S-22, subject to the applicable permissions and confidentiality safeguards.
\BlockTitle{Completion instructions}
Complete the bracketed fields and expand tables as needed. Use an explicit missing-data status rather than zero; justify non-applicability. Identify each value as observed, modeled, assumed, projected, verified or authorized. Preserve adverse outcomes and unresolved findings. Record the exact PD, Ex-Ante, Ex-Post and workbook versions. Cross-references do not replace evidence of implementation during this interval.
\BlockTitle{Reference and scope}
Governing basis: NEBIOT Standard 0.8.2 [1]. The PD [2], Ex-Ante [3] and Ex-Post [4] templates and Environmental Accounts [5] retain edition 1.2. Clause, equation, AP and FRM identifiers refer to [1]; report table identifiers are navigation labels, not new requirements. Use the approved project versions and profile. Complete references appear after the crosswalk.\par
PD means Project Description; EA denotes the dated Ex-Ante Estimate; EP denotes the Ex-Post Outcome Report; MR identifies this Monitoring Report. Use the Ex-Post report as an attached or embedded quantitative annex, not an additional quantity. NEBIOT provides contracted technical services; the proponent and implementing entity retain their S-25 duties. The OVV controls independent assessment. S-26 applies to exceptional NEBIOT interests.
\clearpage
{\fontsize{23}{27}\selectfont\color{NEBGreen}Document control\par}
\begin{Notice}Complete one report for the stated monitoring interval. Preserve the governing PD, the frozen forecast and the reported account as separately identified records. This report supports monitoring submission; it is not an independent verification opinion or a unit-issuance decision.\end{Notice}
\BlockTitle{MR-CTRL / Report record}
{\sffamily\fontsize{8.6}{10.7}\selectfont\begin{longtable}{@{}P{82mm}P{82mm}@{}}
\toprule
\textbf{Field} & \textbf{Record}\\\midrule\endhead
\rule{0pt}{8mm}Template / governing standard & MONITORING / 1.3 / EN; NEBIOT 0.8.2\\\hline
\rule{0pt}{8mm}Project name and identifier & [Complete]\\\hline
\rule{0pt}{8mm}Monitoring report number and project revision & [Complete]\\\hline
\rule{0pt}{8mm}Monitoring start and end / interval convention & [Complete]\\\hline
\rule{0pt}{8mm}Evidence cut-off / submission date & [Complete]\\\hline
\rule{0pt}{8mm}PD / EA / EP / workbook versions & [Complete]\\\hline
\rule{0pt}{8mm}Prepared by / reviewed by / submitted by & [Complete]\\\hline
\rule{0pt}{8mm}Distribution, access and restrictions & [Complete]\\\hline
\end{longtable}}
\BlockTitle{MR-REV / Report revision history}
{\sffamily\fontsize{8.6}{10.7}\selectfont\begin{longtable}{@{}P{48mm}P{61mm}P{53mm}@{}}
\toprule
\textbf{Version / date} & \textbf{Change and affected scope} & \textbf{Authority / status / prior record}\\\midrule\endhead
\rule{0pt}{8mm}[Complete] & [Complete] & [Complete]\\\hline
\rule{0pt}{8mm}[Complete] & [Complete] & [Complete]\\\hline
\end{longtable}}
\clearpage
\tableofcontents
\ReportSection{01}{Monitoring period and results overview}{P-01; P-02; P-04; P-11; U-01; FRM-02}
Identify the report interval, the evidence cut-off date, the category and spatial scope, and the status of every reported result. This is the G2 monitoring submission. G3 verification and G4 authorization are separate decisions; an ex-post result is not automatically verified.\par
\Prompt{Summarize implementation, environmental performance, material events, unresolved findings and the scope submitted for assessment. State monitoring start/end, interval convention, reporting number and evidence cut-off.}
\Answer{16}{[Enter response and exact evidence references.]}
\BlockTitle{MR-01-T01 / Period-result summary — one compatible account}
\par\begingroup\sffamily\fontsize{8.4}{10.5}\selectfont\color{NEBGreen}
\textbf{MR-01-T01 / Calculation notation}\
{\fontsize{10}{12}\selectfont\[\widehat Q_t^z=N_{\mathrm{BL},t}^z-N_{\mathrm{PJ},t}^z-L_t^z,\qquad z\in\{\mathrm{EA},\mathrm{EP}\}\]}
EA denotes the dated forecast; EP denotes the reporting-period account, not automatic verification. All terms are signed period masses on matched support. The quantity is before uncertainty, previous liabilities, reserve, and denomination.\par
\textbf{Basis:} M04.5; F-R01.\quad\textbf{Calculation location:} ExAnte!F:I; ExPost!F:I.\par\endgroup
{\sffamily\fontsize{8.6}{10.7}\selectfont\begin{longtable}{@{}P{48mm}P{61mm}P{53mm}@{}}
\toprule
\textbf{Quantity / symbol} & \textbf{Value and unit} & \textbf{Source cell, period and status}\\\midrule\endhead
\rule{0pt}{8mm}Net baseline emissions | $N_{\mathrm{BL},t}$ & [Complete] & [Complete]\\\hline
\rule{0pt}{8mm}Net project emissions | $N_{\mathrm{PJ},t}$ & [Complete] & [Complete]\\\hline
\rule{0pt}{8mm}Attributable leakage | $L_t$ & [Complete] & [Complete]\\\hline
\rule{0pt}{8mm}Signed result | $\widehat Q_t$ & [Complete] & [Complete]\\\hline
\rule{0pt}{8mm}Conservative result | $Q_t^{\mathrm{low}}$ & [Complete] & [Complete]\\\hline
\rule{0pt}{8mm}Candidate mass | $J_t^{\mathrm{cand}}$ & [Complete] & [Complete]\\\hline
\rule{0pt}{8mm}NVEU candidate, only with category/denomination | $n_t^{\mathrm{candidate}}$ & [Complete] & [Complete]\\\hline
\end{longtable}}
Report each additional environmental or social attribute separately in Section 22. Do not aggregate different attributes or add previous periods to the current-period result. Explain every restriction on the summary in the detailed account.\par
\ReportSection{02}{Controlled PD, forecast and outcome records}{S-18; CYC-05; CYC-06; P-01; P-04; FRM-01; FRM-02}
Use the Project Description (PD) as the design reference, the dated Ex-Ante Estimate as the forecast reference, and the Ex-Post Outcome Report and workbook as the quantitative annex. Reference exact approved or submitted versions and their actual status. A citation to the PD does not replace period-specific evidence.\par
\BlockTitle{MR-02-T01 / Controlled document chain}
{\sffamily\fontsize{8.6}{10.7}\selectfont\begin{longtable}{@{}P{48mm}P{61mm}P{53mm}@{}}
\toprule
\textbf{Record} & \textbf{Version, date and approval/status reference} & \textbf{Exact section, file/fingerprint and access}\\\midrule\endhead
\rule{0pt}{8mm}Governing PD and monitoring plan & [Complete] & [Complete]\\\hline
\rule{0pt}{8mm}Frozen ex-ante forecast & [Complete] & [Complete]\\\hline
\rule{0pt}{8mm}Ex-post report / calculation annex & [Complete] & [Complete]\\\hline
\rule{0pt}{8mm}Environmental Accounts workbook & [Complete] & [Complete]\\\hline
\rule{0pt}{8mm}FRM-05 profile and applicable decisions & [Complete] & [Complete]\\\hline
\rule{0pt}{8mm}Previous MR, opinion and decision & [Complete] & [Complete]\\\hline
\rule{0pt}{8mm}Migration reconciliation, if applicable & [Complete] & [Complete]\\\hline
\end{longtable}}
\Prompt{Identify what is incorporated by exact reference, what is attached, and what changed during the interval. Select one authoritative quantitative annex: an attached Ex-Post report or an embedded copy with equivalent scope. Summaries in this report reproduce that account and do not form a second account.}
\Answer{16}{[Enter response and exact evidence references.]}
\ReportSection{03}{Responsible parties, rights and access}{S-04; S-05; S-21; S-22; S-25; S-26; AO-02; AO-03; FRM-07}
Confirm the actual proponent, implementing entity, rights holders, NEBIOT technical-services lead and data custodians for this interval. Technical monitoring does not make NEBIOT the owner or implementer. Record exceptional NEBIOT interests and the applicable S-26 controls.\par
\BlockTitle{MR-03-T01 / Responsibilities during the interval}
{\sffamily\fontsize{8.6}{10.7}\selectfont\begin{longtable}{@{}P{48mm}P{61mm}P{53mm}@{}}
\toprule
\textbf{Role / entity} & \textbf{Mandate, task and accountable lead} & \textbf{Evidence, effective dates and conflicts}\\\midrule\endhead
\rule{0pt}{8mm}Proponent / implementing entity & [Complete] & [Complete]\\\hline
\rule{0pt}{8mm}Rights holders / representatives & [Complete] & [Complete]\\\hline
\rule{0pt}{8mm}NEBIOT technical services & [Complete] & [Complete]\\\hline
\rule{0pt}{8mm}Field / laboratory / data custody & [Complete] & [Complete]\\\hline
\rule{0pt}{8mm}Internal reviewer / OVV interface & [Complete] & [Complete]\\\hline
\end{longtable}}
\Prompt{Record changes to tenure, mandates, implementation authority, data-use permission, consent, access and continuing obligations. Cite evidence and effective dates; identify unresolved restrictions and the affected scope. Keep restricted personal information in controlled annexes.}
\Answer{16}{[Enter response and exact evidence references.]}
\ReportSection{04}{Implementation against the PD}{CYC-05; S-25; P-02; AO-06; AO-07; AO-08}
Report what the implementing entity actually delivered during the monitoring interval against the PD work program. Distinguish intervention completion, intermediate responses and environmental results. Expenditure, training or a completed activity does not by itself quantify mitigation.\par
\BlockTitle{MR-04-T01 / Commitment-to-evidence matrix}
{\sffamily\fontsize{8.6}{10.7}\selectfont\begin{longtable}{@{}P{48mm}P{61mm}P{53mm}@{}}
\toprule
\textbf{PD activity / indicator and target period} & \textbf{Actual delivery, dates and implementing party} & \textbf{Evidence, deviation and response}\\\midrule\endhead
\rule{0pt}{8mm}[Complete] & [Complete] & [Complete]\\\hline
\rule{0pt}{8mm}[Complete] & [Complete] & [Complete]\\\hline
\rule{0pt}{8mm}[Complete] & [Complete] & [Complete]\\\hline
\rule{0pt}{8mm}[Complete] & [Complete] & [Complete]\\\hline
\end{longtable}}
\Prompt{Explain delays, nonimplementation, maintenance, capacity, safety, access and resource constraints. Link material deviations to the account, additionality assessment and corrective actions; retain favorable and adverse evidence.}
\Answer{16}{[Enter response and exact evidence references.]}
\BlockTitle{MR-04-T02 / Performance-record definitions}
\par\begingroup\sffamily\fontsize{8.4}{10.5}\selectfont\color{NEBGreen}
\textbf{MR-04-T02 / Calculation notation}\
{\fontsize{10}{12}\selectfont\[x_{\ell}\;[u_{\ell}]\]}
x\_l is the identified observation, input, or output; u\_l is its declared unit. Record the exact equation and source file, worksheet, cell/range, period, and status. This notation does not prescribe an estimator.\par
\textbf{Basis:} F-01; T-01.\quad\textbf{Calculation location:} Evidence!A:J; Components!F:L.\par\endgroup
{\sffamily\fontsize{8.6}{10.7}\selectfont\begin{longtable}{@{}P{48mm}P{61mm}P{53mm}@{}}
\toprule
\textbf{Indicator / symbol / unit} & \textbf{Planned value, observed value and interval} & \textbf{Counting or calculation method / evidence}\\\midrule\endhead
\rule{0pt}{8mm}[Complete] & [Complete] & [Complete]\\\hline
\rule{0pt}{8mm}[Complete] & [Complete] & [Complete]\\\hline
\end{longtable}}
\ReportSection{05}{Monitoring plan: intended versus performed}{P-02; P-03; P-04; T-01; FRM-02}
Complete one linked record for every monitored variable required by the PD and applicable modules, including social and ecosystem indicators. Record actual dates and methods rather than repeating only the planned frequency.\par
\BlockTitle{MR-05-T01 / Variable and planned protocol}
\par\begingroup\sffamily\fontsize{8.4}{10.5}\selectfont\color{NEBGreen}
\textbf{MR-05-T01 / Calculation notation}\
{\fontsize{10}{12}\selectfont\[x_{\ell}\;[u_{\ell}]\]}
x\_l is the identified observation, input, or output; u\_l is its declared unit. Record the exact equation and source file, worksheet, cell/range, period, and status. This notation does not prescribe an estimator.\par
\textbf{Basis:} F-01; T-01.\quad\textbf{Calculation location:} Evidence!A:J; Components!F:L.\par\endgroup
{\sffamily\fontsize{8.6}{10.7}\selectfont\begin{longtable}{@{}P{48mm}P{61mm}P{53mm}@{}}
\toprule
\textbf{Variable / symbol / unit and support} & \textbf{PD method, frequency and acceptance criterion} & \textbf{Responsible person, calibration and output equation}\\\midrule\endhead
\rule{0pt}{8mm}[Complete] & [Complete] & [Complete]\\\hline
\rule{0pt}{8mm}[Complete] & [Complete] & [Complete]\\\hline
\rule{0pt}{8mm}[Complete] & [Complete] & [Complete]\\\hline
\end{longtable}}
\BlockTitle{MR-05-T02 / Execution and evidence continuation}
{\sffamily\fontsize{8.6}{10.7}\selectfont\begin{longtable}{@{}P{48mm}P{61mm}P{53mm}@{}}
\toprule
\textbf{Same variable / record ID} & \textbf{Actual dates, method and observation coverage} & \textbf{Raw record, quality outcome and deviation ID}\\\midrule\endhead
\rule{0pt}{8mm}[Complete] & [Complete] & [Complete]\\\hline
\rule{0pt}{8mm}[Complete] & [Complete] & [Complete]\\\hline
\rule{0pt}{8mm}[Complete] & [Complete] & [Complete]\\\hline
\end{longtable}}
\Prompt{Explain omitted campaigns, incomplete coverage and unmeasured variables. State the approved substitution or unresolved status and consequence for the affected result. A blank field is not a zero observation.}
\Answer{16}{[Enter response and exact evidence references.]}
\ReportSection{06}{Spatial scope and monitoring maps}{S-03; S-06; AO-09; P-04; MIG-13; MIG-14; FRM-04/M}
Identify separately the project boundary, accounting strata, reference domain, leakage assessment area and actual support of each result. Include the observation dates, map version, coordinate reference system, legend and disclosure restrictions.\par
\BlockTitle{Monitoring map and versioned spatial reference}
\ProjectMapBox
\Prompt{Map ID / observation date / boundary version / source / legend / authorized-access annex.}
\BlockTitle{MR-06-T01 / Spatial-input and boundary record}
\par\begingroup\sffamily\fontsize{8.4}{10.5}\selectfont\color{NEBGreen}
\textbf{MR-06-T01 / Calculation notation}\
{\fontsize{10}{12}\selectfont\[A_{i,t}=\frac{n_{i,t}a_{\mathrm{cell}}}{k_A},\qquad k_A=10^4\ \mathrm{m^2\,ha^{-1}}\]}
n is a pixel count; a\_cell is pixel area in m2; A is area in ha. The PIXELS route uses the displayed geometric conversion. Direct area is a documented input. Net class change is not gross deforestation.\par
\textbf{Basis:} M01; MIG-13; MIG-14.\quad\textbf{Calculation location:} Spatial!F:I; Profile!B15.\par\endgroup
{\sffamily\fontsize{8.6}{10.7}\selectfont\begin{longtable}{@{}P{48mm}P{61mm}P{53mm}@{}}
\toprule
\textbf{Domain / layer / version} & \textbf{Area and unit; observation dates; source support} & \textbf{Change from PD, comparability and evidence}\\\midrule\endhead
\rule{0pt}{8mm}[Complete] & [Complete] & [Complete]\\\hline
\rule{0pt}{8mm}[Complete] & [Complete] & [Complete]\\\hline
\rule{0pt}{8mm}[Complete] & [Complete] & [Complete]\\\hline
\end{longtable}}
\Prompt{Identify affected boundaries, unmatched dates, unavailable observations and any controlled correction. Map detail or resampling does not establish independent measurement accuracy; retain the processing record and the applicable assessment.}
\Answer{16}{[Enter response and exact evidence references.]}
\ReportSection{07}{NEB001: vegetation identity and change}{T-02; T-03; T-06; T-07; T-21; T-22; M01.4; M01.5; F-F10; MIG-14}
Report the actual classification, CCA calibration and fuzzy-model versions, parameter/rule records, independent reference observations and temporal-comparability checks where applicable. Distinguish vegetation identity from an administrative forest definition and fuzzy membership from probability.\par
\BlockTitle{MR-07-T01 / Reference sampling and area estimates}
\par\begingroup\sffamily\fontsize{8.4}{10.5}\selectfont\color{NEBGreen}
\textbf{MR-07-T01 / Calculation notation}\
{\fontsize{10}{12}\selectfont\[W_i=\frac{A_i}{A},\quad \widehat p_{ij}=\frac{n_{ij}}{n_i},\quad \widehat A_j=A\sum_i W_i\widehat p_{ij}\]}
Use the sampling design and domain required by M01.4. W\_i is the mapped stratum weight, not an estimated weight by default; n\_ij and n\_i are reference counts. This upstream estimator is not fitted by the account workbook.\par
\textbf{Basis:} M01.4; M01.5; T-03.\quad\textbf{Calculation location:} Evidence!G; Spatial!I; qualified upstream calculation.\par\endgroup
{\sffamily\fontsize{8.6}{10.7}\selectfont\begin{longtable}{@{}P{48mm}P{61mm}P{53mm}@{}}
\toprule
\textbf{Map stratum / reference class / interval} & \textbf{Reference counts, weights and estimated area} & \textbf{Sampling design, uncertainty and calculation file}\\\midrule\endhead
\rule{0pt}{8mm}[Complete] & [Complete] & [Complete]\\\hline
\rule{0pt}{8mm}[Complete] & [Complete] & [Complete]\\\hline
\rule{0pt}{8mm}[Complete] & [Complete] & [Complete]\\\hline
\end{longtable}}
\BlockTitle{MR-07-T02 / Gross transitions and their interpretation}
\par\begingroup\sffamily\fontsize{8.4}{10.5}\selectfont\color{NEBGreen}
\textbf{MR-07-T02 / Calculation notation}\
{\fontsize{10}{12}\selectfont\[\Delta A_i=A_i(t_1)-A_i(t_0),\qquad \bar a_i=\frac{\Delta A_i}{t_1-t_0}\]}
The difference is a signed net class-area change, not a gross transition matrix. The average is an allocation over the stated duration, not an annual observation. A gross transition needs its own measured or modeled origin/destination evidence.\par
\textbf{Basis:} M01; MIG-14; reporting identity.\quad\textbf{Calculation location:} Spatial!I:J; Allocation!C:I.\par\endgroup
{\sffamily\fontsize{8.6}{10.7}\selectfont\begin{longtable}{@{}P{48mm}P{61mm}P{53mm}@{}}
\toprule
\textbf{Origin / destination / observation interval} & \textbf{Observed or estimated area (ha) and source ID} & \textbf{Deforestation, degradation, recovery or other; evidence}\\\midrule\endhead
\rule{0pt}{8mm}[Complete] & [Complete] & [Complete]\\\hline
\rule{0pt}{8mm}[Complete] & [Complete] & [Complete]\\\hline
\rule{0pt}{8mm}[Complete] & [Complete] & [Complete]\\\hline
\end{longtable}}
A net class-area difference is not a gross transition estimate. Establish prior state, persistence and physical change; absence of a mapped loss event does not establish absence of degradation. Keep event dates and any temporal allocation distinct.\par
\ReportSection{08}{NEB002: material, carbon and model evidence}{T-04; T-05; T-06; T-07; T-11; T-12; T-13; T-15; M02.1; M02.2; FRM-08}
Identify the A, B or C evidence route actually used for each stratum and pool. Supply new measurements or the qualified reused records required by that route; do not imply a new reconstruction or laboratory campaign where evidence was reused.\par
\BlockTitle{MR-08-T01 / Material-to-carbon record}
\par\begingroup\sffamily\fontsize{8.4}{10.5}\selectfont\color{NEBGreen}
\textbf{MR-08-T01 / Calculation notation}\
{\fontsize{10}{12}\selectfont\[C_j=10^{-3}\sum_o V_{jo}\rho^b_{jo}f^C_{jo}\]}
{\fontsize{10}{12}\selectfont\[c_p=\frac{\sum_{j\in p}C_j+C_{p,\mathrm{other}}}{a_p}\]}
V is green material volume in m3; rho\_b is oven-dry kg per green m3; f\_C is dry-mass carbon fraction. C\_j is Mg C; c\_p is Mg C/ha. The plot denominator a\_p is in ha. Pools must be disjoint and moisture conventions matched.\par
\textbf{Basis:} M02.1; M02.2; F-M01.\quad\textbf{Calculation location:} Evidence!G; Factors!D:G; upstream organ and plot records.\par\endgroup
{\sffamily\fontsize{8.6}{10.7}\selectfont\begin{longtable}{@{}P{48mm}P{61mm}P{53mm}@{}}
\toprule
\textbf{Individual / plot / organ or pool / date} & \textbf{Volume, density, carbon fraction and plot area} & \textbf{Carbon result, uncertainty and evidence location}\\\midrule\endhead
\rule{0pt}{8mm}[Complete] & [Complete] & [Complete]\\\hline
\rule{0pt}{8mm}[Complete] & [Complete] & [Complete]\\\hline
\rule{0pt}{8mm}[Complete] & [Complete] & [Complete]\\\hline
\end{longtable}}
\BlockTitle{MR-08-T02 / Reused estimator and applicability}
{\sffamily\fontsize{8.6}{10.7}\selectfont\begin{longtable}{@{}P{48mm}P{61mm}P{53mm}@{}}
\toprule
\textbf{Route / model or map / actual version} & \textbf{Input domain, validation and transfer checks} & \textbf{Accepted use, restrictions and raw-source reference}\\\midrule\endhead
\rule{0pt}{8mm}[Complete] & [Complete] & [Complete]\\\hline
\rule{0pt}{8mm}[Complete] & [Complete] & [Complete]\\\hline
\rule{0pt}{8mm}[Complete] & [Complete] & [Complete]\\\hline
\end{longtable}}
\Prompt{Describe growth, mortality, recruitment, damage and relevant changes since the PD. Retain failed tests and actual model records. Distinguish a correction of standing stock from a change through time; no universal carbon multiplier is inferred.}
\Answer{16}{[Enter response and exact evidence references.]}
\ReportSection{09}{Quality control, custody and missing observations}{P-03; P-04; T-04; T-06; F-01; A-05}
Trace field records, instruments, laboratory batches, calibration, processing, manual adjustments and archived outputs. Identify who collected, checked and accepted each dataset. Preserve originals and superseded versions.\par
\BlockTitle{MR-09-T01 / Instrument, batch and data acceptance}
{\sffamily\fontsize{8.6}{10.7}\selectfont\begin{longtable}{@{}P{48mm}P{61mm}P{53mm}@{}}
\toprule
\textbf{Dataset / instrument / batch / custodian} & \textbf{Calibration and quality criterion; actual result} & \textbf{Acceptance, limitation or corrective-action ID}\\\midrule\endhead
\rule{0pt}{8mm}[Complete] & [Complete] & [Complete]\\\hline
\rule{0pt}{8mm}[Complete] & [Complete] & [Complete]\\\hline
\rule{0pt}{8mm}[Complete] & [Complete] & [Complete]\\\hline
\end{longtable}}
\BlockTitle{MR-09-T02 / Missing-data and substitution register}
\par\begingroup\sffamily\fontsize{8.4}{10.5}\selectfont\color{NEBGreen}
\textbf{MR-09-T02 / Calculation notation}\
{\fontsize{10}{12}\selectfont\[x_{\ell}\;[u_{\ell}]\]}
x\_l is the identified observation, input, or output; u\_l is its declared unit. Record the exact equation and source file, worksheet, cell/range, period, and status. This notation does not prescribe an estimator.\par
\textbf{Basis:} F-01; T-01.\quad\textbf{Calculation location:} Evidence!A:J; Components!F:L.\par\endgroup
{\sffamily\fontsize{8.6}{10.7}\selectfont\begin{longtable}{@{}P{48mm}P{61mm}P{53mm}@{}}
\toprule
\textbf{Variable / period / missing support} & \textbf{Replacement or bound; unit; uncertainty} & \textbf{Approval, effect on matched accounts and remaining gap}\\\midrule\endhead
\rule{0pt}{8mm}[Complete] & [Complete] & [Complete]\\\hline
\rule{0pt}{8mm}[Complete] & [Complete] & [Complete]\\\hline
\rule{0pt}{8mm}[Complete] & [Complete] & [Complete]\\\hline
\end{longtable}}
\Prompt{Explain how detection limits, inaccessible areas, clouds, lost samples or monitoring failures affect the result. A non-detection is interpreted through the coverage and detection capability of the method, not automatically as zero.}
\Answer{16}{[Enter response and exact evidence references.]}
\ReportSection{10}{Deviations, corrections and parameter changes}{S-18; CYC-06; P-03; P-13; P-14; FRM-05}
Separate a deviation in implementation, a data correction, a change to the method, and a proposed change not yet authorized. Preserve the original PD and frozen forecast, with a dated link to each approved amendment.\par
\BlockTitle{MR-10-T01 / Deviation and decision register}
{\sffamily\fontsize{8.6}{10.7}\selectfont\begin{longtable}{@{}P{48mm}P{61mm}P{53mm}@{}}
\toprule
\textbf{PD provision / planned approach / event date} & \textbf{Actual approach, reason and affected variables} & \textbf{Approval, effective period, finding and remedy}\\\midrule\endhead
\rule{0pt}{8mm}[Complete] & [Complete] & [Complete]\\\hline
\rule{0pt}{8mm}[Complete] & [Complete] & [Complete]\\\hline
\rule{0pt}{8mm}[Complete] & [Complete] & [Complete]\\\hline
\end{longtable}}
\BlockTitle{MR-10-T02 / Parameter change and account effect}
\par\begingroup\sffamily\fontsize{8.4}{10.5}\selectfont\color{NEBGreen}
\textbf{MR-10-T02 / Calculation notation}\
{\fontsize{10}{12}\selectfont\[x_{\ell}\;[u_{\ell}]\]}
x\_l is the identified observation, input, or output; u\_l is its declared unit. Record the exact equation and source file, worksheet, cell/range, period, and status. This notation does not prescribe an estimator.\par
\textbf{Basis:} F-01; T-01.\quad\textbf{Calculation location:} Evidence!A:J; Components!F:L.\par\endgroup
{\sffamily\fontsize{8.6}{10.7}\selectfont\begin{longtable}{@{}P{48mm}P{61mm}P{53mm}@{}}
\toprule
\textbf{Symbol / unit / FRM-05 field} & \textbf{Old and new values, versions and evidence} & \textbf{Baseline, project, leakage and net effects}\\\midrule\endhead
\rule{0pt}{8mm}[Complete] & [Complete] & [Complete]\\\hline
\rule{0pt}{8mm}[Complete] & [Complete] & [Complete]\\\hline
\rule{0pt}{8mm}[Complete] & [Complete] & [Complete]\\\hline
\end{longtable}}
\Prompt{Record AP-01 through AP-12 by applicability, including the exact unchanged-profile reference where retained. State whether revalidation or a restriction is required by the governing rules. Prior validation does not close a newly identified error.}
\Answer{16}{[Enter response and exact evidence references.]}
\ReportSection{11}{Counterfactual for the monitored interval}{S-07; S-08; S-18; AO-05; F-B01; F-B02; P-13; MIG-05}
Identify the approved baseline applicable to this interval, including reference region, drivers, pools, gases and allocation. Distinguish observed conditions from the counterfactual; the baseline does not become an observation because the account is ex post.\par
\BlockTitle{MR-11-T01 / Retained baseline components}
\par\begingroup\sffamily\fontsize{8.4}{10.5}\selectfont\color{NEBGreen}
\textbf{MR-11-T01 / Calculation notation}\
{\fontsize{10}{12}\selectfont\[N_{\mathrm{BL},t}=\sum_{\ell\in\mathcal I_{\mathrm{BL},t}}q_{\ell,\mathrm{BL},t}\]}
Baseline inputs belong to the approved counterfactual and matched period. Do not enter a subtotal that already deducts project emissions as N\_BL and deduct those emissions again. Define each driver, pool, and gas and keep the original cell expression.\par
\textbf{Basis:} M04.4; M04.5; P-13.\quad\textbf{Calculation location:} Components!D=BL; ExAnte!F; ExPost!F.\par\endgroup
{\sffamily\fontsize{8.6}{10.7}\selectfont\begin{longtable}{@{}P{48mm}P{61mm}P{53mm}@{}}
\toprule
\textbf{Driver / pool / gas / period} & \textbf{Input, factor and net component mass} & \textbf{Equation, raw-source cell and approved basis}\\\midrule\endhead
\rule{0pt}{8mm}[Complete] & [Complete] & [Complete]\\\hline
\rule{0pt}{8mm}[Complete] & [Complete] & [Complete]\\\hline
\rule{0pt}{8mm}[Complete] & [Complete] & [Complete]\\\hline
\rule{0pt}{8mm}[Complete] & [Complete] & [Complete]\\\hline
\end{longtable}}
\Prompt{Explain whether the original counterfactual is retained or changed, with the governing decision and matched effect on project and leakage calculations. A subtotal already net of project emissions is not a gross baseline input. Record applicable national reference or nesting obligations through the dated profile.}
\Answer{16}{[Enter response and exact evidence references.]}
\ReportSection{12}{Actual project emissions and removals}{M04.4; M04.5; F-C01; F-C02; T-08; AP-07; P-02}
Report observed project activity and supported conversions, including implementation emissions and any removals within the approved scope. Identify sources as included, excluded under a stated rule, or unresolved. Do not replace observations with an assumed ex-ante effectiveness.\par
\BlockTitle{MR-12-T01 / Gas and source account}
\par\begingroup\sffamily\fontsize{8.4}{10.5}\selectfont\color{NEBGreen}
\textbf{MR-12-T01 / Calculation notation}\
{\fontsize{10}{12}\selectfont\[N_{s,t}=E_{\mathrm{CO_2},s,t}-R_{\mathrm{CO_2},s,t}+\sum_{g\ne\mathrm{CO_2}}\mathrm{GWP}_gE_{g,s,t}\]}
s is BL or PJ; E is emissions and R is CO2 removals. Record gas-specific masses and the adopted characterization source and horizon. Do not characterize an already characterized CO2e quantity again. The workbook imports these qualified component calculations.\par
\textbf{Basis:} M04.4; AP-07.\quad\textbf{Calculation location:} Components!F:L; Profile!B17.\par\endgroup
{\sffamily\fontsize{8.6}{10.7}\selectfont\begin{longtable}{@{}P{48mm}P{61mm}P{53mm}@{}}
\toprule
\textbf{Source / sink / gas / period} & \textbf{Raw amount, unit, factor and net contribution} & \textbf{GWP source/horizon, conversion and evidence}\\\midrule\endhead
\rule{0pt}{8mm}[Complete] & [Complete] & [Complete]\\\hline
\rule{0pt}{8mm}[Complete] & [Complete] & [Complete]\\\hline
\rule{0pt}{8mm}[Complete] & [Complete] & [Complete]\\\hline
\rule{0pt}{8mm}[Complete] & [Complete] & [Complete]\\\hline
\end{longtable}}
\BlockTitle{MR-12-T02 / Component calculation and exclusions}
\par\begingroup\sffamily\fontsize{8.4}{10.5}\selectfont\color{NEBGreen}
\textbf{MR-12-T02 / Calculation notation}\
{\fontsize{10}{12}\selectfont\[q_{\ell,s,t}=x_{\ell,s,t}f_{\ell,s,t}\]}
{\fontsize{10}{12}\selectfont\[N_{s,t}=\sum_{\ell\in\mathcal I_{s,t}}q_{\ell,s,t}\]}
x and f have reciprocal units yielding t CO2e per period. I is the set of included, evidenced, disjoint components for the same basis, period and account. For precomputed mass, f=1 must be explicitly entered with its meaning. BL and PJ are net emissions; LK is attributable leakage.\par
\textbf{Basis:} M04.4; M04.5; T-08; arithmetic implementation.\quad\textbf{Calculation location:} Components!G:I; ExAnte!F:H; ExPost!F:H.\par\endgroup
{\sffamily\fontsize{8.6}{10.7}\selectfont\begin{longtable}{@{}P{48mm}P{61mm}P{53mm}@{}}
\toprule
\textbf{Component / symbol / included status} & \textbf{Activity value, factor and output unit} & \textbf{Source cell, zero/exclusion basis and review}\\\midrule\endhead
\rule{0pt}{8mm}[Complete] & [Complete] & [Complete]\\\hline
\rule{0pt}{8mm}[Complete] & [Complete] & [Complete]\\\hline
\rule{0pt}{8mm}[Complete] & [Complete] & [Complete]\\\hline
\end{longtable}}
\ReportSection{13}{Offsite activity and attributable leakage}{S-09; F-L01; T-08; T-16; P-03; P-04}
Distinguish observed offsite loss from displacement attributable to the project. Use matched support for the with-project and without-project offsite accounts, and identify all attribution adjustments and restrictions.\par
\BlockTitle{MR-13-T01 / Leakage evidence and matched account}
\par\begingroup\sffamily\fontsize{8.4}{10.5}\selectfont\color{NEBGreen}
\textbf{MR-13-T01 / Calculation notation}\
{\fontsize{10}{12}\selectfont\[\widehat L_t=N_{\mathrm{PLKA,with},t}-N_{\mathrm{PLKA,without},t}\]}
The two accounts have matched geography, activities, gases and intervals. Any source ratio remains an explicit numerator/denominator, not an assumed percentage. A signed diagnostic does not automatically authorize an extra benefit; the applicable attribution profile governs.\par
\textbf{Basis:} F-L01; T-08.\quad\textbf{Calculation location:} Components!D=LK; ExAnte!H; ExPost!H.\par\endgroup
{\sffamily\fontsize{8.6}{10.7}\selectfont\begin{longtable}{@{}P{48mm}P{61mm}P{53mm}@{}}
\toprule
\textbf{Activity / offsite domain / interval} & \textbf{With-project and without-project amounts; unit} & \textbf{Attributed leakage, adjustment and source evidence}\\\midrule\endhead
\rule{0pt}{8mm}[Complete] & [Complete] & [Complete]\\\hline
\rule{0pt}{8mm}[Complete] & [Complete] & [Complete]\\\hline
\rule{0pt}{8mm}[Complete] & [Complete] & [Complete]\\\hline
\rule{0pt}{8mm}[Complete] & [Complete] & [Complete]\\\hline
\end{longtable}}
\Prompt{Record actual surveys, reference evidence, drivers, included pools and gases, temporal allocation, remaining uncertainty and the governing rule for signed diagnostics. Explain every zero or exclusion. Keep components already included in project emissions out of a second leakage deduction.}
\Answer{16}{[Enter response and exact evidence references.]}
\ReportSection{14}{Stocks, delayed releases and cohort continuity}{M04.1; M04.2; M04.3; F-C03a; F-C03b; F-C04; T-08; MIG-07}
Reconcile the same physical loss through one consistent stock-change or event-release pathway. Retain cohorts originating before the current reporting interval and before migration where they still release material. Distinguish biomass, carbon, CO2 and CO2e.\par
\BlockTitle{MR-14-T01 / Cohort release and remaining material}
\par\begingroup\sffamily\fontsize{8.4}{10.5}\selectfont\color{NEBGreen}
\textbf{MR-14-T01 / Calculation notation}\
{\fontsize{10}{12}\selectfont\[E^{\mathrm{CO_2}}_{i,t}=\frac{44}{12}\sum_{\tau\le t}\sum_o A_{i,\tau}\ell_{io,\tau}k_{io,\tau}(t)\]}
A is newly affected ha, ell is allocated Mg C/ha, and k is the fraction released as CO2 in interval t. The sum of fractions cannot exceed its defined allocation. Cohorts and pools are disjoint; keep remaining material and pre-period cohorts. Use exact intervals.\par
\textbf{Basis:} M04.2; F-C03a; F-C03b; F-C04.\quad\textbf{Calculation location:} Upstream cohort schedule; Components!F:L.\par\endgroup
{\sffamily\fontsize{8.6}{10.7}\selectfont\begin{longtable}{@{}P{48mm}P{61mm}P{53mm}@{}}
\toprule
\textbf{Stratum / pool / event interval} & \textbf{Allocated carbon, release kernel and current release} & \textbf{Remaining allocation, balance and evidence}\\\midrule\endhead
\rule{0pt}{8mm}[Complete] & [Complete] & [Complete]\\\hline
\rule{0pt}{8mm}[Complete] & [Complete] & [Complete]\\\hline
\rule{0pt}{8mm}[Complete] & [Complete] & [Complete]\\\hline
\rule{0pt}{8mm}[Complete] & [Complete] & [Complete]\\\hline
\end{longtable}}
\Prompt{Explain pool transfers, harvested material, decay, burning and gas characterization where applicable. Check that previously recognized releases are not counted twice and that the sequence does not release more than its documented allocation. A revised carbon stock is not itself a newly observed removal.}
\Answer{16}{[Enter response and exact evidence references.]}
\ReportSection{15}{Signed ex-post account and reporting totals}{M04.5; F-R01; F-R03; P-04; P-13; FRM-02}
\BlockTitle{MR-15-T01 / Period and stratum account}
\par\begingroup\sffamily\fontsize{8.4}{10.5}\selectfont\color{NEBGreen}
\textbf{MR-15-T01 / Calculation notation}\
{\fontsize{10}{12}\selectfont\[\widehat Q_t^z=N_{\mathrm{BL},t}^z-N_{\mathrm{PJ},t}^z-L_t^z,\qquad z\in\{\mathrm{EA},\mathrm{EP}\}\]}
EA denotes the dated forecast; EP denotes the reporting-period account, not automatic verification. All terms are signed period masses on matched support. The quantity is before uncertainty, previous liabilities, reserve, and denomination.\par
\textbf{Basis:} M04.5; F-R01.\quad\textbf{Calculation location:} ExAnte!F:I; ExPost!F:I.\par\endgroup
{\sffamily\fontsize{8.6}{10.7}\selectfont\begin{longtable}{@{}P{48mm}P{61mm}P{53mm}@{}}
\toprule
\textbf{Period / stratum / category} & \textbf{Baseline, project and leakage masses (t CO2e)} & \textbf{Signed result and exact workbook row}\\\midrule\endhead
\rule{0pt}{8mm}[Complete] & [Complete] & [Complete]\\\hline
\rule{0pt}{8mm}[Complete] & [Complete] & [Complete]\\\hline
\rule{0pt}{8mm}[Complete] & [Complete] & [Complete]\\\hline
\rule{0pt}{8mm}[Complete] & [Complete] & [Complete]\\\hline
\end{longtable}}
\Prompt{Identify disjoint included components and the method for aggregating them. Reconcile the reported total to the authoritative Ex-Post annex and workbook. Retain negative contributions, excluded scope and unresolved inputs before any nonnegative eligible portion.}
\Answer{16}{[Enter response and exact evidence references.]}
\BlockTitle{MR-15-T02 / Current interval versus cumulative history}
\par\begingroup\sffamily\fontsize{8.4}{10.5}\selectfont\color{NEBGreen}
\textbf{MR-15-T02 / Calculation notation}\
{\fontsize{10}{12}\selectfont\[x_{\ell}\;[u_{\ell}]\]}
x\_l is the identified observation, input, or output; u\_l is its declared unit. Record the exact equation and source file, worksheet, cell/range, period, and status. This notation does not prescribe an estimator.\par
\textbf{Basis:} F-01; T-01.\quad\textbf{Calculation location:} Evidence!A:J; Components!F:L.\par\endgroup
{\sffamily\fontsize{8.6}{10.7}\selectfont\begin{longtable}{@{}P{48mm}P{61mm}P{53mm}@{}}
\toprule
\textbf{Quantity / basis / interval} & \textbf{Current interval; previous recognized periods; unit} & \textbf{Cumulative method, overlap control and evidence}\\\midrule\endhead
\rule{0pt}{8mm}[Complete] & [Complete] & [Complete]\\\hline
\rule{0pt}{8mm}[Complete] & [Complete] & [Complete]\\\hline
\rule{0pt}{8mm}[Complete] & [Complete] & [Complete]\\\hline
\end{longtable}}
Cumulative results are supporting history, not an additional claim. Do not sum the total with its components, annual totals with their monthly allocations, or this report with its attached Ex-Post report.\par
\ReportSection{16}{Uncertainty, conservativeness and materiality}{M05.1; M05.2; M05.3; F-U01; F-A01; F-A02; T-14; T-16; P-13; AP-08; AP-10}
\BlockTitle{MR-16-T01 / Uncertainty components and dependence}
\par\begingroup\sffamily\fontsize{8.4}{10.5}\selectfont\color{NEBGreen}
\textbf{MR-16-T01 / Calculation notation}\
{\fontsize{10}{12}\selectfont\[u_Q^2=u_B^2+u_P^2+u_L^2-2\operatorname{Cov}(B,P)-2\operatorname{Cov}(B,L)+2\operatorname{Cov}(P,L)\]}
B, P and L in this variance identity denote baseline, project and leakage estimators, not reserve B\_t. Standard uncertainties are in t CO2e; variances and covariances are squared units. Retain shared-source dependence and a valid covariance matrix.\par
\textbf{Basis:} F-U01; M05.1; M05.2.\quad\textbf{Calculation location:} Qualified uncertainty calculation; ExAnte!J:L; ExPost!J:L.\par\endgroup
{\sffamily\fontsize{8.6}{10.7}\selectfont\begin{longtable}{@{}P{48mm}P{61mm}P{53mm}@{}}
\toprule
\textbf{Estimator / component / unit} & \textbf{Uncertainty, covariance and supported bias treatment} & \textbf{Evidence, propagation and unresolved restriction}\\\midrule\endhead
\rule{0pt}{8mm}[Complete] & [Complete] & [Complete]\\\hline
\rule{0pt}{8mm}[Complete] & [Complete] & [Complete]\\\hline
\rule{0pt}{8mm}[Complete] & [Complete] & [Complete]\\\hline
\end{longtable}}
\BlockTitle{MR-16-T02 / Conservative result and separate materiality decision}
\par\begingroup\sffamily\fontsize{8.4}{10.5}\selectfont\color{NEBGreen}
\textbf{MR-16-T02 / Calculation notation}\
{\fontsize{10}{12}\selectfont\[Q_t^{\mathrm{low}}=\widehat Q_t-k_{1-\alpha}u_{Q_t}\]}
{\fontsize{10}{12}\selectfont\[Q_t^{\mathrm{low}}=\widehat Q_t-\delta_t\quad(\mathrm{DEDUCTION});\qquad Q_t^{\mathrm{low}}=q_t^{\mathrm{bound}}\quad(\mathrm{BOUND})\]}
delta is the justified deduction (DEDUCTION route); BOUND supplies the supported lower value directly. The second displayed line describes the workbook interface, not a new approved estimator. Record the upstream uncertainty or evidence-bound procedure and decision.\par
\textbf{Basis:} M05.3; F-A02; AP-08.\quad\textbf{Calculation location:} ExAnte!J:L; ExPost!J:L; ExPost!O.\par\endgroup
{\sffamily\fontsize{8.6}{10.7}\selectfont\begin{longtable}{@{}P{48mm}P{61mm}P{53mm}@{}}
\toprule
\textbf{Period / selected method / AP reference} & \textbf{Inputs, supported bound or deduction and result} & \textbf{Materiality criterion, identified errors and disposition}\\\midrule\endhead
\rule{0pt}{8mm}[Complete] & [Complete] & [Complete]\\\hline
\rule{0pt}{8mm}[Complete] & [Complete] & [Complete]\\\hline
\rule{0pt}{8mm}[Complete] & [Complete] & [Complete]\\\hline
\end{longtable}}
No confidence level, deduction or materiality threshold is prefilled. Explain the joint evidence behind the bound and avoid deducting twice for the same covered component. Correct identified arithmetic errors; a difference below a threshold does not turn an error into a valid input.\par
\ReportSection{17}{Risk, reserve and disturbance response}{T-17; T-18; T-19; T-20; F-A03; F-A04; F-A05; P-03; P-10; AP-09; FRM-09}
\BlockTitle{MR-17-T01 / Residual-risk and protection account}
\par\begingroup\sffamily\fontsize{8.4}{10.5}\selectfont\color{NEBGreen}
\textbf{MR-17-T01 / Calculation notation}\
{\fontsize{10}{12}\selectfont\[B^{\mathrm{req}}=\inf\{\ell:\Pr(L_H\le\ell\mid\mathcal D_v)\ge\gamma\}\]}
{\fontsize{10}{12}\selectfont\[B^{\mathrm{add}}=\max(0,B^{\mathrm{req}}-B^{\mathrm{avail}})\]}
This is conditional F-A04, not a default probability model. Record the adopted horizon, protection level and loss distribution, or the separately justified stress route. Distinguish total required protection, available compatible reserve and new contribution. No rate is prefilled.\par
\textbf{Basis:} F-A04; F-A05; T-17–T-20; AP-09.\quad\textbf{Calculation location:} Qualified risk schedule; Settlement!J; FRM-09.\par\endgroup
{\sffamily\fontsize{8.6}{10.7}\selectfont\begin{longtable}{@{}P{48mm}P{61mm}P{53mm}@{}}
\toprule
\textbf{Hazard / compatible exposure / horizon} & \textbf{Implemented controls, loss evidence and protection rule} & \textbf{Required, available and additional reserve; evidence}\\\midrule\endhead
\rule{0pt}{8mm}[Complete] & [Complete] & [Complete]\\\hline
\rule{0pt}{8mm}[Complete] & [Complete] & [Complete]\\\hline
\rule{0pt}{8mm}[Complete] & [Complete] & [Complete]\\\hline
\end{longtable}}
\BlockTitle{MR-17-T02 / Disturbances, response and current status}
\par\begingroup\sffamily\fontsize{8.4}{10.5}\selectfont\color{NEBGreen}
\textbf{MR-17-T02 / Calculation notation}\
{\fontsize{10}{12}\selectfont\[x_{\ell}\;[u_{\ell}]\]}
x\_l is the identified observation, input, or output; u\_l is its declared unit. Record the exact equation and source file, worksheet, cell/range, period, and status. This notation does not prescribe an estimator.\par
\textbf{Basis:} F-01; T-01.\quad\textbf{Calculation location:} Evidence!A:J; Components!F:L.\par\endgroup
{\sffamily\fontsize{8.6}{10.7}\selectfont\begin{longtable}{@{}P{48mm}P{61mm}P{53mm}@{}}
\toprule
\textbf{Event, date and affected support} & \textbf{Measured or estimated impact and unit} & \textbf{Immediate action, evidence, liability and follow-up}\\\midrule\endhead
\rule{0pt}{8mm}[Complete] & [Complete] & [Complete]\\\hline
\rule{0pt}{8mm}[Complete] & [Complete] & [Complete]\\\hline
\rule{0pt}{8mm}[Complete] & [Complete] & [Complete]\\\hline
\end{longtable}}
Keep required protection, funded reserve, committed reserve and executed cancellation distinct. Record dependence and common shocks where relevant. A stronger assessment does not itself create more environmental quantity. Do not identify every disturbance or negative bound as a verified reversal without the applicable assessment.\par
\ReportSection{18}{Settlement, negative results and continuity}{F-03; M05.4; M05.5; F-R03; F-R05a; F-R05b; FRM-02/C; FRM-09; AP-09}
\BlockTitle{MR-18-T01 / Candidate allocation and executed posting}
\par\begingroup\sffamily\fontsize{8.4}{10.5}\selectfont\color{NEBGreen}
\textbf{MR-18-T01 / Calculation notation}\
{\fontsize{10}{12}\selectfont\[Q_t^{\mathrm{eligible}}=\max(0,Q_t^{\mathrm{low}}),\quad D_t^{\mathrm{cand}}=\min(D_t,Q_t^{\mathrm{eligible}})\]}
{\fontsize{10}{12}\selectfont\[Q_t^{\mathrm{base,cand}}=Q_t^{\mathrm{eligible}}-D_t^{\mathrm{cand}},\quad J_t^{\mathrm{cand}}=Q_t^{\mathrm{base,cand}}-B_t\]}
cand distinguishes a worksheet proposal from an executed posting. D\_set is actual evidenced settlement; D\_rem = D - D\_set. Do not reduce posted liability by D\_cand. A larger required reserve than the candidate base is a shortfall, not funded protection. No rounding is implied.\par
\textbf{Basis:} M05.4; M05.5; F-R03.\quad\textbf{Calculation location:} Settlement!B:L; ExAnte!M:O (forecast only).\par\endgroup
{\sffamily\fontsize{8.6}{10.7}\selectfont\begin{longtable}{@{}P{48mm}P{61mm}P{53mm}@{}}
\toprule
\textbf{Opening liability and eligible result} & \textbf{Candidate versus executed settlement; evidence} & \textbf{Post-liability base, reserve and remaining mass}\\\midrule\endhead
\rule{0pt}{8mm}[Complete] & [Complete] & [Complete]\\\hline
\rule{0pt}{8mm}[Complete] & [Complete] & [Complete]\\\hline
\rule{0pt}{8mm}[Complete] & [Complete] & [Complete]\\\hline
\end{longtable}}
\BlockTitle{MR-18-T02 / Negative-result continuation}
\par\begingroup\sffamily\fontsize{8.4}{10.5}\selectfont\color{NEBGreen}
\textbf{MR-18-T02 / Calculation notation}\
{\fontsize{10}{12}\selectfont\[N_t^-=\max(0,-Q_t^{\mathrm{low}}),\quad \Delta D_t^-=N_t^- -O_t-C_t^-\]}
{\fontsize{10}{12}\selectfont\[D_{t+1}=D_t^{\mathrm{rem}}+\Delta D_t^-\]}
FULL\_BOUND only when expressly adopted. O is a documented previously recognized portion of the same deficit; C\_minus is executed reserve cancellation for a disjoint portion. Require O+C\_minus <= N\_minus. A negative uncertainty bound is not automatically a physical reversal.\par
\textbf{Basis:} F-R05a; F-R05b; F-03.\quad\textbf{Calculation location:} CarryForward!B:J; Settlement!I.\par\endgroup
{\sffamily\fontsize{8.6}{10.7}\selectfont\begin{longtable}{@{}P{48mm}P{61mm}P{53mm}@{}}
\toprule
\textbf{Period / rule / signed conservative result} & \textbf{New deficit, prior recognized overlap and canceled reserve} & \textbf{Unsettled increment, closing liability and next opening}\\\midrule\endhead
\rule{0pt}{8mm}[Complete] & [Complete] & [Complete]\\\hline
\rule{0pt}{8mm}[Complete] & [Complete] & [Complete]\\\hline
\rule{0pt}{8mm}[Complete] & [Complete] & [Complete]\\\hline
\end{longtable}}
Apply the full-bound recurrence only when expressly adopted in FRM-05. Verify nonnegative compatible balances and no double recognition of the same deficit. Earmarking reserve is not executed cancellation. Retain the actual settlement evidence and reconcile the next opening balance; a zero eligible quantity does not extinguish an obligation.\par
\ReportSection{19}{Comparison with the frozen ex-ante forecast}{S-18; P-13; P-14; M04.5; F-R01}
\BlockTitle{MR-19-T01 / Matched forecast and reported account}
\par\begingroup\sffamily\fontsize{8.4}{10.5}\selectfont\color{NEBGreen}
\textbf{MR-19-T01 / Calculation notation}\
{\fontsize{10}{12}\selectfont\[\Delta Q_t=Q_t^{\mathrm{EP}}-Q_t^{\mathrm{EA}}\]}
{\fontsize{10}{12}\selectfont\[\Delta Q_t=\Delta N_{\mathrm{BL},t}-\Delta N_{\mathrm{PJ},t}-\Delta L_t\]}
All deltas are EP minus the frozen EA value on the same support before uncertainty and reserve. r = deltaQ - (deltaBL - deltaPJ - deltaL) is an arithmetic residual. A forecast difference is not itself mitigation, causation, or a materiality finding.\par
\textbf{Basis:} M04.5; P-13; comparison identity.\quad\textbf{Calculation location:} Comparison!D:J.\par\endgroup
{\sffamily\fontsize{8.6}{10.7}\selectfont\begin{longtable}{@{}P{48mm}P{61mm}P{53mm}@{}}
\toprule
\textbf{Period and matched scope} & \textbf{EA and EP values for baseline, project and leakage} & \textbf{Net difference, arithmetic closure and explanation}\\\midrule\endhead
\rule{0pt}{8mm}[Complete] & [Complete] & [Complete]\\\hline
\rule{0pt}{8mm}[Complete] & [Complete] & [Complete]\\\hline
\rule{0pt}{8mm}[Complete] & [Complete] & [Complete]\\\hline
\end{longtable}}
\Prompt{Separate actual performance, scope/time harmonization, new observations, correction of factors, method revisions, uncertainty, reserves and liability effects. Keep the original forecast alongside any harmonized comparison. Agreement is not required where a change is justified; identify each difference rather than forcing it to zero.}
\Answer{16}{[Enter response and exact evidence references.]}
\BlockTitle{MR-19-T02 / Subperiod allocation, only if used}
\par\begingroup\sffamily\fontsize{8.4}{10.5}\selectfont\color{NEBGreen}
\textbf{MR-19-T02 / Calculation notation}\
{\fontsize{10}{12}\selectfont\[Q_{t,m}=w_{t,m}Q_t,\quad w_{t,m}\ge0,\quad \sum_m w_{t,m}=1\]}
{\fontsize{10}{12}\selectfont\[r_t^{\mathrm{alloc}}=\sum_m Q_{t,m}-Q_t\]}
m is a child reporting interval, not automatically a measured month. The same parent and basis must be used for all weights. Do not add a parent total to its allocated children. Retain negative signs and explain unequal-duration allocations.\par
\textbf{Basis:} F-R01; allocation identity.\quad\textbf{Calculation location:} Allocation!C:I.\par\endgroup
{\sffamily\fontsize{8.6}{10.7}\selectfont\begin{longtable}{@{}P{48mm}P{61mm}P{53mm}@{}}
\toprule
\textbf{Parent quantity / interval / status} & \textbf{Weights, child intervals and allocation method} & \textbf{Allocated total, closure and source record}\\\midrule\endhead
\rule{0pt}{8mm}[Complete] & [Complete] & [Complete]\\\hline
\rule{0pt}{8mm}[Complete] & [Complete] & [Complete]\\\hline
\end{longtable}}
\ReportSection{20}{Continuing additionality and overlapping initiatives}{S-08; AO-01; AO-02; AO-05; AO-06; AO-07; AO-08; AO-09; AO-10; AO-11; AO-12; FRM-02/AO; FRM-04/AO}
Reassess the applicable additionality premises against implementation, changed threats, legal duties, common practice and financing where required by the selected test. Report existing and concurrent interventions affecting the counterfactual; a new applicant or map does not imply an untreated territory.\par
\BlockTitle{MR-20-T01 / Additionality continuation record}
{\sffamily\fontsize{8.6}{10.7}\selectfont\begin{longtable}{@{}P{48mm}P{61mm}P{53mm}@{}}
\toprule
\textbf{PD premise / intervention / review trigger} & \textbf{Actual action, intermediate response and outcome evidence} & \textbf{Competing explanation, current conclusion and restriction}\\\midrule\endhead
\rule{0pt}{8mm}[Complete] & [Complete] & [Complete]\\\hline
\rule{0pt}{8mm}[Complete] & [Complete] & [Complete]\\\hline
\rule{0pt}{8mm}[Complete] & [Complete] & [Complete]\\\hline
\rule{0pt}{8mm}[Complete] & [Complete] & [Complete]\\\hline
\end{longtable}}
\BlockTitle{MR-20-T02 / Overlap and outcome reconciliation}
{\sffamily\fontsize{8.6}{10.7}\selectfont\begin{longtable}{@{}P{48mm}P{61mm}P{53mm}@{}}
\toprule
\textbf{Intersecting initiative / period / support} & \textbf{Authority, activities and potential shared result} & \textbf{Unique scope, resolution, hold and evidence}\\\midrule\endhead
\rule{0pt}{8mm}[Complete] & [Complete] & [Complete]\\\hline
\rule{0pt}{8mm}[Complete] & [Complete] & [Complete]\\\hline
\rule{0pt}{8mm}[Complete] & [Complete] & [Complete]\\\hline
\end{longtable}}
Geographic intersection alone does not settle authority, additionality or duplication. Do not assign the full common environmental effect to every initiative. Record unresolved material findings and their effect on the submitted scope.\par
\ReportSection{21}{Participation, safeguards and benefit-sharing}{S-12; S-13; S-17; AO-03; P-12; A-13; A-14; A-15}
Document actual engagement, consent where applicable, governance decisions, worker and community safeguards, benefit-sharing and access to complaints during the interval. Distinguish planned commitments from executed actions and payments.\par
\BlockTitle{MR-21-T01 / Commitments and observed performance}
{\sffamily\fontsize{8.6}{10.7}\selectfont\begin{longtable}{@{}P{48mm}P{61mm}P{53mm}@{}}
\toprule
\textbf{Commitment / group / period} & \textbf{Executed measure, evidence and distribution basis} & \textbf{Adverse effect, outstanding duty and response}\\\midrule\endhead
\rule{0pt}{8mm}[Complete] & [Complete] & [Complete]\\\hline
\rule{0pt}{8mm}[Complete] & [Complete] & [Complete]\\\hline
\rule{0pt}{8mm}[Complete] & [Complete] & [Complete]\\\hline
\rule{0pt}{8mm}[Complete] & [Complete] & [Complete]\\\hline
\end{longtable}}
\BlockTitle{MR-21-T02 / Complaint and remedy follow-up}
{\sffamily\fontsize{8.6}{10.7}\selectfont\begin{longtable}{@{}P{48mm}P{61mm}P{53mm}@{}}
\toprule
\textbf{Restricted case ID / date / issue} & \textbf{Channel, approved deadline and independent handling} & \textbf{Response, remedy evidence and closure status}\\\midrule\endhead
\rule{0pt}{8mm}[Complete] & [Complete] & [Complete]\\\hline
\rule{0pt}{8mm}[Complete] & [Complete] & [Complete]\\\hline
\rule{0pt}{8mm}[Complete] & [Complete] & [Complete]\\\hline
\end{longtable}}
\Prompt{Explain language and nondigital access, protections against retaliation, consent changes, vulnerable-group impacts and unresolved matters. The public report may redact protected identities, but authorized assessment needs the full evidence. Do not treat absence of recorded complaints as proof that access was effective.}
\Answer{16}{[Enter response and exact evidence references.]}
\ReportSection{22}{Ecosystem, biodiversity and social outcomes}{M06.1; T-10; T-22; S-02; P-11; U-01; U-02}
Report only the indicators and methods supported by the PD and applicable profile. Keep biodiversity, ecosystem condition and social outcomes in their own units. Record targets, observed period values and cumulative values separately; a lifetime target is not an achieved monitoring-period result.\par
\BlockTitle{MR-22-T01 / Indicator-to-evidence record}
\par\begingroup\sffamily\fontsize{8.4}{10.5}\selectfont\color{NEBGreen}
\textbf{MR-22-T01 / Calculation notation}\
{\fontsize{10}{12}\selectfont\[x_{\ell}\;[u_{\ell}]\]}
x\_l is the identified observation, input, or output; u\_l is its declared unit. Record the exact equation and source file, worksheet, cell/range, period, and status. This notation does not prescribe an estimator.\par
\textbf{Basis:} F-01; T-01.\quad\textbf{Calculation location:} Evidence!A:J; Components!F:L.\par\endgroup
{\sffamily\fontsize{8.6}{10.7}\selectfont\begin{longtable}{@{}P{48mm}P{61mm}P{53mm}@{}}
\toprule
\textbf{Indicator / definition / symbol / unit} & \textbf{Reference, period result and uncertainty / limitation} & \textbf{Method, observation dates, source and attribution}\\\midrule\endhead
\rule{0pt}{8mm}[Complete] & [Complete] & [Complete]\\\hline
\rule{0pt}{8mm}[Complete] & [Complete] & [Complete]\\\hline
\rule{0pt}{8mm}[Complete] & [Complete] & [Complete]\\\hline
\rule{0pt}{8mm}[Complete] & [Complete] & [Complete]\\\hline
\end{longtable}}
\Prompt{Specify sampling effort, counting rules, repeated observations and any offsite adverse effects. Distinguish unique participants from attendances, jobs from person-time and a species list from demonstrated change, where those indicators apply. Document the PD definition instead of inserting a new universal metric.}
\Answer{16}{[Enter response and exact evidence references.]}
A reported indicator does not become an NVEU without the applicable category and authorization. No conversion between these outcomes and carbon is implied.\par
\ReportSection{23}{Migration and prior claim continuity}{MIG-02; MIG-04; MIG-05; MIG-06; MIG-07; MIG-08; MIG-09; MIG-12; MIG-17; AO-11}
Complete when migration, reuse of historical accounts or prior claims affect the interval. Preserve the originating method, dates, rights, findings, issued/retired/canceled history and liabilities. Attach the controlled Migration Reconciliation; do not restart the baseline or obligation clock by changing the report.\par
\BlockTitle{MR-23-T01 / Original, harmonized and NEBIOT account}
\par\begingroup\sffamily\fontsize{8.4}{10.5}\selectfont\color{NEBGreen}
\textbf{MR-23-T01 / Calculation notation}\
{\fontsize{10}{12}\selectfont\[\Delta q_{\mathrm{basis}}=q^{\mathrm{harm}}-q^{\mathrm{src}},\quad \Delta q_{\mathrm{method}}=q^{\mathrm{NEB}}-q^{\mathrm{harm}}\]}
{\fontsize{10}{12}\selectfont\[q^{\mathrm{NEB}}-q^{\mathrm{src}}=\Delta q_{\mathrm{basis}}+\Delta q_{\mathrm{method}}\]}
src retains the original result; harm is the original method on the agreed comparable basis; NEB is the receiving-method result on that basis. This telescoping identity explains differences, not eligibility. Do not force the methodological difference to zero.\par
\textbf{Basis:} MIG-05; P-13; reconciliation identity.\quad\textbf{Calculation location:} Migration!E:K.\par\endgroup
{\sffamily\fontsize{8.6}{10.7}\selectfont\begin{longtable}{@{}P{48mm}P{61mm}P{53mm}@{}}
\toprule
\textbf{Quantity / unit / original scope} & \textbf{Original, harmonized and receiving-method values} & \textbf{Basis change, method change and evidence}\\\midrule\endhead
\rule{0pt}{8mm}[Complete] & [Complete] & [Complete]\\\hline
\rule{0pt}{8mm}[Complete] & [Complete] & [Complete]\\\hline
\rule{0pt}{8mm}[Complete] & [Complete] & [Complete]\\\hline
\end{longtable}}
\BlockTitle{MR-23-T02 / Prior claims and continuing obligations}
{\sffamily\fontsize{8.6}{10.7}\selectfont\begin{longtable}{@{}P{48mm}P{61mm}P{53mm}@{}}
\toprule
\textbf{Outcome / period / program record} & \textbf{Other external carbon units and original status} & \textbf{Overlap, remaining liability and receiving disposition}\\\midrule\endhead
\rule{0pt}{8mm}[Complete] & [Complete] & [Complete]\\\hline
\rule{0pt}{8mm}[Complete] & [Complete] & [Complete]\\\hline
\rule{0pt}{8mm}[Complete] & [Complete] & [Complete]\\\hline
\end{longtable}}
Cancellation alone does not authorize replacement issuance. Retaining a qualified historical map is not a waiver of evidence or uniqueness. Use the same underlying result only once.\par
\ReportSection{24}{NVEU request and claim boundaries}{U-01; U-02; U-03; U-05; U-07; F-R04; P-09; P-11; FRM-06}
Identify whether the submission requests only assessment of environmental results or also a later NVEU decision. State the specific category, denomination, precision, permitted claim and governing approval where these exist. Leave an unapproved category pending and report physical quantities instead.\par
\BlockTitle{MR-24-T01 / Conditional category calculation}
\par\begingroup\sffamily\fontsize{8.4}{10.5}\selectfont\color{NEBGreen}
\textbf{MR-24-T01 / Calculation notation}\
{\fontsize{10}{12}\selectfont\[n_t^{\mathrm{candidate}}=\frac{J_t^{\mathrm{cand}}}{d_u},\qquad d_u>0,\quad J_t^{\mathrm{cand}}\ge0\]}
d\_u is t CO2e per NVEU of the specified category. The quotient is a candidate count, not issuance or an equivalence with other external carbon units. No value for d\_u or whole-unit rule is assumed. Missing category authorization blocks the unit count.\par
\textbf{Basis:} F-R04; U-01.\quad\textbf{Calculation location:} Results!F:H; Profile!B7:B10.\par\endgroup
{\sffamily\fontsize{8.6}{10.7}\selectfont\begin{longtable}{@{}P{48mm}P{61mm}P{53mm}@{}}
\toprule
\textbf{Category / period / support} & \textbf{Candidate mass, denomination and NVEU count} & \textbf{Approval, rounding/residual rule and exact cell}\\\midrule\endhead
\rule{0pt}{8mm}[Complete] & [Complete] & [Complete]\\\hline
\rule{0pt}{8mm}[Complete] & [Complete] & [Complete]\\\hline
\rule{0pt}{8mm}[Complete] & [Complete] & [Complete]\\\hline
\end{longtable}}
\BlockTitle{MR-24-T02 / Reported, verified and authorized status}
\par\begingroup\sffamily\fontsize{8.4}{10.5}\selectfont\color{NEBGreen}
\textbf{MR-24-T02 / Calculation notation}\
{\fontsize{10}{12}\selectfont\[x_{\ell}\;[u_{\ell}]\]}
x\_l is the identified observation, input, or output; u\_l is its declared unit. Record the exact equation and source file, worksheet, cell/range, period, and status. This notation does not prescribe an estimator.\par
\textbf{Basis:} F-01; T-01.\quad\textbf{Calculation location:} Evidence!A:J; Components!F:L.\par\endgroup
{\sffamily\fontsize{8.6}{10.7}\selectfont\begin{longtable}{@{}P{48mm}P{61mm}P{53mm}@{}}
\toprule
\textbf{Same category / quantity / unit} & \textbf{Reported / independently verified / authorized amounts} & \textbf{Opinion, decision and status-history reference}\\\midrule\endhead
\rule{0pt}{8mm}[Complete] & [Complete] & [Complete]\\\hline
\rule{0pt}{8mm}[Complete] & [Complete] & [Complete]\\\hline
\rule{0pt}{8mm}[Complete] & [Complete] & [Complete]\\\hline
\end{longtable}}
\Prompt{Reconcile the request with prior recognition, use, retirement, cancellation and remaining liabilities. Separate current-request quantities from history, reserves and denomination residuals. No status changes merely because a quantity appears in the monitoring report. A pending opinion or decision remains pending.}
\Answer{16}{[Enter response and exact evidence references.]}
\ReportSection{25}{Findings, submission and separate decisions}{S-15; S-16; S-22; P-05; P-06; P-09; P-14; A-12; A-13; FRM-03; FRM-06}
\BlockTitle{MR-25-T01 / Findings and independent closure}
{\sffamily\fontsize{8.6}{10.7}\selectfont\begin{longtable}{@{}P{48mm}P{61mm}P{53mm}@{}}
\toprule
\textbf{Requirement / finding / affected scope} & \textbf{Proponent response, corrected evidence and date} & \textbf{Uninvolved assessment, closure and remaining condition}\\\midrule\endhead
\rule{0pt}{8mm}[Complete] & [Complete] & [Complete]\\\hline
\rule{0pt}{8mm}[Complete] & [Complete] & [Complete]\\\hline
\rule{0pt}{8mm}[Complete] & [Complete] & [Complete]\\\hline
\rule{0pt}{8mm}[Complete] & [Complete] & [Complete]\\\hline
\end{longtable}}
Retain each original finding, response and independent disposition. A completed response or future action does not close a material issue in the current claim. Identify the approved materiality and assurance scope, actual OVV credentials and access arrangements by reference.\par
\BlockTitle{MR-25-T02 / Submission and decision references}
{\sffamily\fontsize{8.6}{10.7}\selectfont\begin{longtable}{@{}P{48mm}P{61mm}P{53mm}@{}}
\toprule
\textbf{Role / record} & \textbf{Name, date, signature or signed-record reference} & \textbf{Scope and status}\\\midrule\endhead
\rule{0pt}{8mm}Prepared: NEBIOT technical-services lead & [Complete] & [Complete]\\\hline
\rule{0pt}{8mm}Submitted: accountable proponent & [Complete] & [Complete]\\\hline
\rule{0pt}{8mm}Internal review: uninvolved competent reviewer & [Complete] & [Complete]\\\hline
\rule{0pt}{8mm}OVV opinion / program decision, when issued & [Complete] & [Complete]\\\hline
\end{longtable}}
The last row references a separately issued record; it is not a form for the preparer to sign on behalf of the OVV or program authority. Record pending assessment explicitly. Keep reported, verified and authorized quantities separate.\par
\ReportSection{26}{Evidence handover and next-period actions}{P-04; P-10; P-12; AO-08; CYC-06; S-18; FRM-02; FRM-06}
\BlockTitle{MR-26-T01 / Evidence and annex index}
{\sffamily\fontsize{8.6}{10.7}\selectfont\begin{longtable}{@{}P{48mm}P{61mm}P{53mm}@{}}
\toprule
\textbf{Evidence ID / document / version / fingerprint} & \textbf{Exact record, section, worksheet or cell range} & \textbf{Custodian, access, retention approval and report use}\\\midrule\endhead
\rule{0pt}{8mm}[Complete] & [Complete] & [Complete]\\\hline
\rule{0pt}{8mm}[Complete] & [Complete] & [Complete]\\\hline
\rule{0pt}{8mm}[Complete] & [Complete] & [Complete]\\\hline
\rule{0pt}{8mm}[Complete] & [Complete] & [Complete]\\\hline
\rule{0pt}{8mm}[Complete] & [Complete] & [Complete]\\\hline
\end{longtable}}
\BlockTitle{MR-26-T02 / Continuing actions and unresolved conditions}
{\sffamily\fontsize{8.6}{10.7}\selectfont\begin{longtable}{@{}P{48mm}P{61mm}P{53mm}@{}}
\toprule
\textbf{Action / original finding / trigger} & \textbf{Responsible party, due date or approved frequency} & \textbf{Evidence required, affected decision and next report}\\\midrule\endhead
\rule{0pt}{8mm}[Complete] & [Complete] & [Complete]\\\hline
\rule{0pt}{8mm}[Complete] & [Complete] & [Complete]\\\hline
\rule{0pt}{8mm}[Complete] & [Complete] & [Complete]\\\hline
\end{longtable}}
\Prompt{Identify the next monitoring handover, opening balances and cohort records, unchanged or revised plan, outstanding remedies, data-retention arrangements and distribution version. Keep public and controlled evidence indexes linked without exposing protected information. No deadline is assumed by this template.}
\Answer{16}{[Enter response and exact evidence references.]}
\ReportSection{CALC}{Calculation traceability and reviewer recalculation}{F-01; T-01; P-04; P-13; FRM-02; NEB-MTH-002}
Repeat this record for each material calculated quantity or link an equivalent controlled calculation ledger. Reproduce the exact original expression first and record every adaptation separately. Report text, equations and workbook rows must use the same units, period, scope and version. Administrative records do not require invented equations.
\BlockTitle{Notation and interpretation}
$\mathrm{EA}$ / $\mathrm{EP}$: dated forecast / period account; $t$: interval; $i,j$: stratum or class; $o$: organ or pool; $g$: gas; $\tau$: event cohort; $\ell$: calculation row.
\BlockTitle{MR-CALC / Repeatable calculation record}
{\sffamily\fontsize{8.6}{10.7}\selectfont\begin{longtable}{@{}P{82mm}P{82mm}@{}}
\toprule
\textbf{Field} & \textbf{Evidence / value / location}\\\midrule\endhead
\rule{0pt}{8mm}Calculation ID / symbol / intended result & [Complete]\\\hline
\rule{0pt}{8mm}Exact source file / version / fingerprint & [Complete]\\\hline
\rule{0pt}{8mm}Worksheet / cell / original observation ID & [Complete]\\\hline
\rule{0pt}{8mm}Original formula, units and period & [Complete]\\\hline
\rule{0pt}{8mm}Input symbols, values, units and status & [Complete]\\\hline
\rule{0pt}{8mm}NEBIOT equation / original or adapted method & [Complete]\\\hline
\rule{0pt}{8mm}Conversion, inclusion and exclusion rationale & [Complete]\\\hline
\rule{0pt}{8mm}Unrounded intermediate and final results & [Complete]\\\hline
\rule{0pt}{8mm}MR table / Ex-Post table / receiving cell & [Complete]\\\hline
\rule{0pt}{8mm}Arithmetic difference and tolerance settings & [Complete]\\\hline
\rule{0pt}{8mm}Uncertainty and materiality: separate records & [Complete]\\\hline
\rule{0pt}{8mm}Reviewer / evidence / finding / disposition & [Complete]\\\hline
\end{longtable}}
\BlockTitle{Recalculation check}
\[\delta_{\ell}=v_{\ell}^{\mathrm{recalc}}-v_{\ell}^{\mathrm{reference}}\]
\[|\delta_{\ell}|\le\varepsilon_a+\varepsilon_r|v_{\ell}^{\mathrm{reference}}|\]
The absolute tolerance has the result unit and the relative tolerance is dimensionless. Enter both as software comparison settings, not materiality criteria. Compare unrounded values after confirming matched methods and support. A justified NEBIOT methodological difference need not be zero.
\ReportSection{XREF}{PD–Ex-Ante–Ex-Post crosswalk}{CYC-05; P-02; P-04; FRM-01; FRM-02}
The section numbers below refer to template edition 1.2 [2–4], not to every completed project document. Section 02 records the exact project versions and corresponding sections. Follow these references with the evidence and formula records; the crosswalk does not authorize an omitted requirement.
{\sffamily\fontsize{8.2}{10.2}\selectfont\begin{longtable}{@{}P{83mm}P{26mm}P{24mm}P{24mm}@{}}\toprule
MR section & PD [2] & EA [3] & EP [4]\\\midrule\endhead
\hyperref[mr:01]{01 Monitoring period and results overview} & 01; 12 & 01; 07 & 01; 07; 08\\\hline
\hyperref[mr:02]{02 Controlled PD, forecast and outcome records} & 01; 14; 15; 17; A & 01; 11 & 01; 12\\\hline
\hyperref[mr:03]{03 Responsible parties, rights and access} & 02 & 11 & 01; 10\\\hline
\hyperref[mr:04]{04 Implementation against the PD} & 04; 13 & 04; 10 & 04; 10\\\hline
\hyperref[mr:05]{05 Monitoring plan: intended versus performed} & 14; 15 & 02; 08 & 02\\\hline
\hyperref[mr:06]{06 Spatial scope and monitoring maps} & 03; 08 & 02 & 02; 11\\\hline
\hyperref[mr:07]{07 NEB001: vegetation identity and change} & 05; 14 & 02; 06 & 02\\\hline
\hyperref[mr:08]{08 NEB002: material, carbon and model evidence} & 05; 09; 14 & 02; 06 & 02; 06\\\hline
\hyperref[mr:09]{09 Quality control, custody and missing observations} & 05; 14; 17 & 02; 11 & 02; 12\\\hline
\hyperref[mr:10]{10 Deviations, corrections and parameter changes} & 14; 15; 16 & 01; 08; 11 & 02; 03; 09; 12\\\hline
\hyperref[mr:11]{11 Counterfactual for the monitored interval} & 06; 07 & 03 & 03\\\hline
\hyperref[mr:12]{12 Actual project emissions and removals} & 09; 13; 14 & 04; 06 & 04; 06\\\hline
\hyperref[mr:13]{13 Offsite activity and attributable leakage} & 06; 09; 14 & 05 & 05\\\hline
\hyperref[mr:14]{14 Stocks, delayed releases and cohort continuity} & 09; 10 & 06 & 06\\\hline
\hyperref[mr:15]{15 Signed ex-post account and reporting totals} & 09; A; CALC & 07 & 07\\\hline
\hyperref[mr:16]{16 Uncertainty, conservativeness and materiality} & 10; 15 & 08 & 07; CALC\\\hline
\hyperref[mr:17]{17 Risk, reserve and disturbance response} & 10; 13 & 08 & 08\\\hline
\hyperref[mr:18]{18 Settlement, negative results and continuity} & 10; 15 & 07; 08 & 08\\\hline
\hyperref[mr:19]{19 Comparison with the frozen ex-ante forecast} & 09; 10; A & 01; 07; 08 & 09\\\hline
\hyperref[mr:20]{20 Continuing additionality and overlapping initiatives} & 07; 08 & 10 & 10; 11\\\hline
\hyperref[mr:21]{21 Participation, safeguards and benefit-sharing} & 11; 13 & 10 & 10\\\hline
\hyperref[mr:22]{22 Ecosystem, biodiversity and social outcomes} & 04; 11; 12; 14 & 10 & 10\\\hline
\hyperref[mr:23]{23 Migration and prior claim continuity} & 08; A & 09 & 11\\\hline
\hyperref[mr:24]{24 NVEU request and claim boundaries} & 12; 15; 18 & 07; 10 & 08; 12\\\hline
\hyperref[mr:25]{25 Findings, submission and separate decisions} & 16; 18 & 11 & 12\\\hline
\hyperref[mr:26]{26 Evidence handover and next-period actions} & 14; 17; 18 & 11 & 12\\\hline
\end{longtable}}
\ReportSection{REF}{References and completion boundary}{NEBIOT Standard 0.8.2}
[1] NEBIOT. NEBIOT Standard, controlled version 0.8.2. Integrated edition NEB-DOC-000. Applicable core requirements and Annex A; methods NEB-MTH-001/002; project cycle NEB-CYC-001; assurance NEB-PRO-001; NVEU NEB-UNT-001; migration NEB-MIG-001; assessment NEB-AUD-001; controlled forms NEB-FRM-001.\par\vspace{2mm}
[2] NEBIOT. Project Description template. Edition 1.2. Sections 01–18, account-attachment appendix A and CALC appendix. The actual project PD and its authorized changes are identified in Section 02 of this report.\par\vspace{2mm}
[3] NEBIOT. Ex-Ante Estimate template. Edition 1.2. Sections 01–11 and CALC appendix. Retain the dated project forecast used for comparison.\par\vspace{2mm}
[4] NEBIOT. Ex-Post Outcome Report template. Edition 1.2. Sections 01–12 and CALC appendix. Attach or embed the controlled quantitative report for the same interval and scope.\par\vspace{2mm}
[5] NEBIOT. Environmental Accounts workbook and mathematical-notation catalog. Edition 1.2. Record the actual file version and fingerprint. Worksheet and column references in this report identify the unchanged edition 1.2 layout; fill the exact row and cell used.\par\vspace{2mm}
[6] NEBIOT. Migration Reconciliation template. Edition 1.2. Apply to affected historical and current accounts only when migration or prior claims are relevant.\par\vspace{2mm}
\BlockTitle{Completion checks before submission}
Confirm that Section 01 reproduces the same authoritative period account as Sections 15–19 and the attached Ex-Post report; verify units and unrounded totals before formatting.\par
Confirm that each applicable PD monitoring obligation has an actual observation, a justified non-applicability decision, or a clearly identified unresolved gap. Confirm public and controlled evidence references are accessible to their authorized recipients.\par
Confirm prior periods, reserve cancellation, liabilities and continuing cohorts are reconciled once. Leave pending OVV opinions and program decisions pending rather than supplying a prefilled conclusion.\par
Tables are editable records, not automatic calculators. Word, LaTeX and Excel do not synchronize values automatically. Preserve the exact report-to-workbook references and inspect the completed document after tables or fields expand.\par
